\documentclass[11pt]{article}
\usepackage[margin=1in]{geometry}
\usepackage{amsmath, amssymb, amsthm}
\usepackage{mathtools}
\usepackage{hyperref}
\usepackage{algorithm}
\usepackage{algpseudocode}
\usepackage{graphicx}
\usepackage{float}
\usepackage[numbers,sort&compress]{natbib}

\newtheorem{proposition}{Proposition}
\newtheorem{definition}{Definition}

\title{What Algorithms Can Transformers Learn In-Context?\\ Online selection of online algorithms}
\author{Senem I\c{s}\i{}k}

\begin{document}
\maketitle

\tableofcontents
\bigskip

\section{Introduction}
\label{sec:intro}

In-context learning is the phenomenon by which a sequence model, at test time, executes a procedure inferred from the data it sees in context --- potentially noisy input-output pairs only --- without any weight updates. Prior work has demonstrated this for simple function classes such as linear and boolean functions, but not for online algorithms. Yet many real decision problems are solved not by fitting a function but by running such an algorithm: integrating the data so far, updating beliefs, choosing an adaptive strategy, and acting. Can transformers learn this kind of instance adaptability --- recovering and running the right algorithmic strategy for each evolving sequence --- or do they collapse onto a single fixed procedure tuned to the average of the training distribution?

Our testbed is a Bayesian stopping problem with a Normal--Normal conjugate prior, in which the optimal policy is a backward-induction threshold rule indexed by a latent noise scale $\sigma$. We call this setting \emph{online selection of online algorithms}: at test time the transformer sees a sequence with unknown noise level and must, on the fly, recover from the prefix it has observed the threshold rule that corresponds to that noise level.

By varying whether the noise level is fixed or random across training distributions, we engineer a setting that discriminates between two hypotheses about the trained model:
\begin{itemize}
\item \textbf{Learning-algorithm hypothesis.} The transformer infers the noise level in-context and applies the corresponding threshold rule, adapting per instance and generalizing across noise levels.
\item \textbf{Shortcut hypothesis.} The transformer applies a single training-distribution-tuned policy that ignores the noise level, matching the oracle in-distribution but degrading sharply elsewhere.
\end{itemize}
To tell them apart, we analyze the trained network with an interpretability stack --- attention entropy, layerwise probing, residual-stream interventions, activation patching, head ablation, and a logit lens through depth --- asking which intermediate quantities the model computes, where in the network they live, and whether they causally drive the output. The goal is to support the case that transformers can select online algorithms in-context (with the right training), and to explain how.

\section{Setup: A Bayesian Optimal Stopping Problem}
\label{sec:setup}

\subsection{Generative model}

For each sample $i \in \{1, \ldots, N\}$, draw a within-sequence noise scale $\sigma_i$ from a prior, draw a latent mean $\mu_i \sim \mathcal{N}(\mu_0, \tau_0^2)$, then draw a sequence
\[
X_1, X_2, \ldots, X_n \mid \mu_i, \sigma_i \;\overset{\text{i.i.d.}}{\sim}\; \mathcal{N}(\mu_i, \sigma_i^2).
\]
We fix $\mu_0 = 0$, $\tau_0 = 10$, horizon $n = 256$ throughout. The agent observes $X_t$ at each $t \in \{1, \ldots, n\}$ and chooses to \emph{stop} (collect $X_t$) or \emph{continue}; at $t = n$ it must stop. Both $\mu_i$ and (in the unknown-$\sigma$ setting) $\sigma_i$ are latent.

\subsection{Five training distributions}

We instantiate five training distributions, three with $\sigma$ fixed (\emph{static-variance}) and two with $\sigma$ random (\emph{random-variance}):
\begin{itemize}
\item $\mathcal{D}_1$: static, $\sigma = 1$ (data-dominant).
\item $\mathcal{D}_2$: static, $\sigma = 10$ (balanced).
\item $\mathcal{D}_3$: static, $\sigma = 100$ (prior-dominant).
\item $\mathcal{D}_{\mathrm{disc}}$: random, $\sigma \sim \mathrm{Uniform}\{1, 10, 100\}$, each with probability $1/3$.
\item $\mathcal{D}_{\mathrm{logu}}$: random, $\log_{10}\sigma \sim \mathrm{Uniform}(0, 2)$, equivalently $p(\sigma) = 1/(2\sigma\ln 10)$ on $[1, 100]$.
\end{itemize}
Both random-variance priors give equal weight to each decade of $\sigma$ --- the discrete prior trivially, the continuous prior by uniform-in-$\log_{10}\sigma$ construction.

The static-variance distributions $\mathcal{D}_1, \mathcal{D}_2, \mathcal{D}_3$ play a double role: training distributions for the single-regime models, and test regimes for all evaluation (Section~\ref{sec:regimes}). We name single-regime models by their training data (D1\_cv, D1\_act, etc.) when we need to disambiguate.

\subsection{Three test regimes}
\label{sec:regimes}

The static-variance distributions $\mathcal{D}_1, \mathcal{D}_2, \mathcal{D}_3$ are also the three test regimes for all evaluation. They correspond to noise-to-prior-precision ratios $\rho := \sigma^2/\tau_0^2$, with the data weight at $t = 1$ given by $1/(1 + \rho)$:

\begin{center}
\small
\begin{tabular}{l c c c c l}
\hline
test regime & $\sigma$ & $\sigma^2$ & $\rho$ & data weight at $t = 1$ & character \\
\hline
$\mathcal{D}_1$ (data-dominant)    & $1$   & $1$     & $0.01$ & $0.99$ & plug-in is asymptotically Bayes \\
$\mathcal{D}_2$ (balanced)         & $10$  & $100$   & $1$    & $0.50$ & no closed-form is asymptotic \\
$\mathcal{D}_3$ (prior-dominant)   & $100$ & $10000$ & $100$  & $0.01$ & prior-only is asymptotically Bayes \\
\hline
\end{tabular}
\end{center}

\paragraph{Datasets.} Training sequences are sampled fresh from the training distribution at each step (no fixed training set); the meaningful budget is gradient steps $\times$ batch size, specified in Section~\ref{sec:training}. Validation and test sets are pre-generated and held fixed at $N_{\mathrm{val}} = N_{\mathrm{test}} = 10^4$ sequences each (validation drawn from the training distribution, used for checkpoint selection only; test drawn per regime $\mathcal{D}_i$, used for all evaluation). Validation and test sets use disjoint random seeds.

\section{Bayes-optimal Oracle}
\label{sec:oracle}

\subsection{Static-variance setting}
\label{sec:oracle-static}

When the training distribution is one of $\mathcal{D}_1, \mathcal{D}_2, \mathcal{D}_3$ the noise scale $\sigma$ is constant and known; only $\mu$ is latent. The oracle is the standard Normal--Normal stopping DP.

\paragraph{Sufficient statistic and posterior.} The sufficient statistic for $\mu$ given $X_{1:t}$ is the running sum
\[
S_t \;:=\; \sum_{s=1}^t X_s,
\]
with derived sample mean $\bar X_t := S_t / t$ (the maximum-likelihood estimate of $\mu$ at known $\sigma$). The conjugate Normal--Normal posterior depends on $X_{1:t}$ only through $S_t$:
\begin{equation}
\label{eq:mu-static}
\mu \mid X_{1:t} \;\sim\; \mathcal{N}(\mu_t, \tau_t^2),
\qquad
\tau_t^2 = \!\left(\frac{1}{\tau_0^2} + \frac{t}{\sigma^2}\right)^{\!-1},
\qquad
\mu_t = \tau_t^2\!\left(\frac{\mu_0}{\tau_0^2} + \frac{S_t}{\sigma^2}\right).
\end{equation}
$\tau_t^2$ is deterministic in $t$; $\mu_t$ is a deterministic function of $(t, S_t)$. The belief state is $b_t = \mu_t$, one-dimensional.

\paragraph{DP recursion.} The continuation value satisfies $V_t(b_t, X_t) = \max\{X_t, C_t(\mu_t)\}$ for $t < n$ and $V_n = X_n$, with
\begin{equation}
\label{eq:Crec-static}
C_t(\mu_t) \;=\; \mathbb{E}_{X \sim \mathcal{N}(\mu_t,\, \tau_t^2 + \sigma^2)}\!\Big[\max\!\big\{X,\, C_{t+1}(\mathrm{update}(\mu_t, X))\big\}\Big],
\end{equation}
where $\mathrm{update}(\mu, X) = \alpha_t \mu + (1 - \alpha_t) X$ with $\alpha_t := \sigma^2/(\tau_t^2 + \sigma^2)$, and $C_{n-1}(\mu) = \mu$. The Bayes-optimal policy is $\pi^\star_{\mathcal{D}_i, t}(b_t, X_t) = \mathbf{1}[X_t \ge C_t(\mu_t)]$.

\paragraph{Approximate DP.}
\label{sec:adp-static}
Equation~\eqref{eq:Crec-static} has no elementary closed form. We discretize $\mu_t$ on a uniform grid of $K$ points covering $[\mu_0 - 5 s_t,\, \mu_0 + 5 s_t]$ at each stage, where $s_t = \sqrt{\tau_0^2 - \tau_t^2}$ is the marginal standard deviation of $\mu_t$ across sequences. The Gaussian expectation in~\eqref{eq:Crec-static} is approximated by $J$-point Gauss--Hermite quadrature, the standard rule for integrals against $e^{-z^2}$. The nodes $\{z_j\}_{j=1}^J \subset \mathbb{R}$ are the roots of the $J$-th physicists' Hermite polynomial $H_J$, and the weights $\{w_j\}_{j=1}^J > 0$ are uniquely determined by requiring that
\[
\int_{-\infty}^\infty f(z)\, e^{-z^2}\, dz \;\approx\; \sum_{j=1}^J w_j\, f(z_j)
\]
hold with equality whenever $f$ is a polynomial of degree at most $2J - 1$. Both nodes and weights depend only on $J$ and are precomputed (e.g., via \texttt{numpy.polynomial.hermite.hermgauss}). Substituting $X = \mu + \sqrt{2 v_t}\,z$ for $X \sim \mathcal{N}(\mu, v_t)$ with $v_t := \tau_t^2 + \sigma^2$ gives
\begin{equation}
\label{eq:gh}
\mathbb{E}_{X \sim \mathcal{N}(\mu, v_t)}[f(X)] \;\approx\; \frac{1}{\sqrt\pi}\sum_{j=1}^J w_j\, f\!\left(\mu + \sqrt{2 v_t}\, z_j\right).
\end{equation}
Off-grid values of $\widehat C_{t+1}$ are looked up by linear interpolation with clipping at the grid boundary. Algorithm~\ref{alg:adp-static} runs the backward induction at $K = 2048$, $J = 128$.

\begin{algorithm}[H]
\caption{Approximate DP for the static-variance Bayes-optimal threshold (per regime).}
\label{alg:adp-static}
\begin{algorithmic}[1]
\State Precompute $\tau_t^2, v_t := \tau_t^2 + \sigma^2, \alpha_t$ for each $t$; grid $\{m_t^{(k)}\}_{k=1}^K$ at each $t$; Gauss--Hermite nodes $\{(z_j, w_j)\}_{j=1}^J$
\State $\widehat C_{n-1}(m_{n-1}^{(k)}) \gets m_{n-1}^{(k)}$ for $k = 1, \ldots, K$
\For{$t = n-2, n-3, \ldots, 1$}
  \For{$k = 1, \ldots, K$}
    \State $S \gets 0$
    \For{$j = 1, \ldots, J$}
      \State $x_j \gets m_t^{(k)} + \sqrt{2 v_t}\, z_j$
      \State $\mu_j' \gets \alpha_t\, m_t^{(k)} + (1-\alpha_t)\, x_j$
      \State $S \gets S + w_j\, \max\!\big\{x_j,\, \widehat C_{t+1}^{\mathrm{lin}}(\mu_j')\big\}$
    \EndFor
    \State $\widehat C_t(m_t^{(k)}) \gets S/\sqrt\pi$
  \EndFor
\EndFor
\State \Return $\{\widehat C_t\}$
\end{algorithmic}
\end{algorithm}

\paragraph{Convergence check.} Solve at $(K, J) = (2048, 128)$ and at reference $(K, J) = (4096, 256)$; report (i) max absolute disagreement on $\widehat C$ over the grid, and (ii) action-label disagreement on $10^4$ test sequences per regime. Fix the labeling resolution after the operational disagreement falls below $10^{-3}$.

\subsection{Random-variance setting}
\label{sec:oracle-random}

When the training distribution is $\mathcal{D}_{\mathrm{disc}}$ or $\mathcal{D}_{\mathrm{logu}}$, both $\mu$ and $\sigma$ are latent. The oracle is a two-dimensional DP over a joint posterior on $(\mu, \sigma)$.

\paragraph{Sufficient statistics.} The sufficient statistics for the joint posterior are
\begin{equation}
\label{eq:S-Q-defs}
S_t \;:=\; \sum_{s=1}^t X_s,
\qquad
Q_t \;:=\; \sum_{s=1}^t X_s^2.
\end{equation}
Both update online: $S_t = S_{t-1} + X_t$, $Q_t = Q_{t-1} + X_t^2$. Together with $t$, the pair $(S_t, Q_t)$ determines the joint posterior $p(\mu, \sigma \mid X_{1:t})$. $S_t$ alone determines the sample mean $\bar X_t = S_t/t$; jointly $(S_t, Q_t)$ determine the within-sequence sample variance
\[
\hat\sigma_t^2 \;:=\; \frac{1}{t-1}(Q_t - t\bar X_t^2).
\]

\paragraph{Joint posterior.} By Bayes' rule, $p(\mu, \sigma \mid X_{1:t}) \propto p(\sigma)\, p(\mu)\, p(X_{1:t} \mid \mu, \sigma)$. Factor as $p(\sigma \mid X_{1:t}) \cdot p(\mu \mid X_{1:t}, \sigma)$. For any fixed $\sigma$, the conditional posterior $p(\mu \mid X_{1:t}, \sigma)$ is the standard Normal--Normal posterior~\eqref{eq:mu-static}: depends on $X_{1:t}$ only through $S_t$, with different $\sigma$ values giving different conditional means and variances.

For the marginal posterior $p(\sigma \mid X_{1:t})$: the marginal likelihood $p(X_{1:t} \mid \sigma) = \int p(X_{1:t} \mid \mu, \sigma)\, p(\mu)\, d\mu$ has $X_{1:t} \mid \sigma$ multivariate Gaussian with mean $\mathbf{0}$ and covariance $\Sigma_t^{(\sigma)} = \sigma^2 I_t + \tau_0^2 \mathbf{1}\mathbf{1}^\top$. Using the matrix determinant lemma and Sherman--Morrison (Appendix~\ref{app:marginal-derivation}), the log marginal likelihood depends on $X_{1:t}$ only through $(S_t, Q_t)$:
\begin{equation}
\label{eq:marginal-loglik}
\log p(X_{1:t} \mid \sigma) = \mathrm{const}(t) - \frac{t}{2}\log\sigma^2 - \frac{1}{2}\log\!\left(1 + \frac{t\tau_0^2}{\sigma^2}\right) - \frac{1}{2}\!\left[\frac{Q_t}{\sigma^2} - \frac{\tau_0^2 S_t^2/\sigma^4}{1 + t\tau_0^2/\sigma^2}\right].
\end{equation}
For $\mathcal{D}_{\mathrm{disc}}$ the posterior over $\sigma$ is a categorical on three points; for $\mathcal{D}_{\mathrm{logu}}$ a continuous density on $[1, 100]$. Evaluation details are deferred to the Approximate DP paragraph below.

\paragraph{Predictive distribution and DP recursion.} The predictive $p(X_{t+1} \mid X_{1:t})$ is a mixture of Gaussians with weights from $p(\sigma \mid X_{1:t})$:
\begin{equation}
\label{eq:predictive}
p(X_{t+1} \mid X_{1:t}) \;=\; \int p(\sigma \mid X_{1:t}) \cdot \mathcal{N}\!\left(X_{t+1};\, \mu_t^{(\sigma)},\, v_t^{(\sigma)}\right)\, d\sigma,
\qquad v_t^{(\sigma)} := \tau_t^{2,(\sigma)} + \sigma^2,
\end{equation}
with the integral replaced by a sum over $\{1, 10, 100\}$ for $\mathcal{D}_{\mathrm{disc}}$. The belief state is $b_t = (S_t, Q_t)$. Backward induction gives $V_t(b_t, X_t) = \max\{X_t, C_t(b_t)\}$ for $t < n$ with $V_n = X_n$, and
\begin{equation}
\label{eq:Crec-random}
C_t(b_t) \;=\; \mathbb{E}_{X_{t+1} \sim p(\cdot \mid b_t)}\!\big[\max\{X_{t+1}, C_{t+1}(b_{t+1})\}\big],
\end{equation}
expectation over the mixture predictive~\eqref{eq:predictive}, $b_{t+1} = (S_t + X_{t+1}, Q_t + X_{t+1}^2)$, terminal $C_{n-1}(b_{n-1}) = \int p(\sigma \mid X_{1:n-1}) \mu_{n-1}^{(\sigma)}\, d\sigma$. The Bayes-optimal policy is $\pi^\star_{\mathcal{D}, t}(b_t, X_t) = \mathbf{1}[X_t \ge C_t(b_t)]$.

\paragraph{Approximate DP.}
\label{sec:adp-random}
We discretize $b_t$ in transformed coordinates $(\bar X_t, L_t)$ where $L_t := \log\hat\sigma_t^2$. The two axes need different treatment:
\begin{itemize}
\item $\bar X_t = \mu + (1/t)\sum_s \varepsilon_s$ has marginal variance $\tau_0^2 + \sigma^2/t$ across sequences; at fixed $\sigma$, this only converges to $\tau_0^2$ as $t \to \infty$. At small $t$ and large $\sigma$ the marginal SD is much larger than $\tau_0$, so a fixed grid is unsuitable. We instead use a \emph{per-stage} uniform grid $I_{\bar X, t} := [-5\, s_{\bar X, t},\, +5\, s_{\bar X, t}]$ with stage-dependent half-width $s_{\bar X, t} := \sqrt{\tau_0^2 + \sigma_{\max}^2 / t}$, where $\sigma_{\max}$ is the supremum of the prior's support ($\sigma_{\max} = 100$ for both $\mathcal{D}_{\mathrm{disc}}$ and $\mathcal{D}_{\mathrm{logu}}$). At $t = 1$: $s_{\bar X, 1} = \sqrt{\tau_0^2 + \sigma_{\max}^2} \approx 100.5$, grid spans $\approx [-503, 503]$. As $t \to \infty$: $s_{\bar X, t} \to \tau_0 = 10$, grid spans $\approx [-50, 50]$. Mirrors the static ADP's per-stage $\mu_t$ grid (Section~\ref{sec:adp-static}).
\item $L_t$ estimates $\log\sigma^2$, which spans $[\log 1, \log 10^4] = [0, 4\log 10]$. Grid on $I_L := [\log 0.1, \log 10^5]$ to absorb sample-variance dispersion at small $t$. Stage-invariant.
\end{itemize}
Use $K_1 = K_2 = 256$ uniform grid points on each axis, total $K_1 K_2 = 65{,}536$ states per stage. The state translation between $(S_t, Q_t)$ and grid coordinates is $S_t = t \bar X_t$, $Q_t = (t-1) e^{L_t} + t \bar X_t^2$.

For the mixture quadrature, use $M$ components: $M = 3$ for $\mathcal{D}_{\mathrm{disc}}$ at $\sigma \in \{1, 10, 100\}$, or $M = J_\sigma = 64$ Gauss--Legendre nodes on $\log\sigma \in [0, \log 100]$ for $\mathcal{D}_{\mathrm{logu}}$ (Gauss--Legendre is the analog of Gauss--Hermite for integrals against the uniform measure on a finite interval). For each component $\sigma_k$ apply $J = 64$-point Gauss--Hermite~\eqref{eq:gh}; sum the per-component results weighted by $\tilde w_k := p(\sigma_k \mid X_{1:t}) \cdot \omega_k / Z_t$, where $\omega_k$ is the Gauss--Legendre weight (or $1$ for the discrete case) and $Z_t$ normalizes. Off-grid lookup uses 2D bilinear interpolation in $(\bar X_{t+1}, L_{t+1})$.

\begin{algorithm}[H]
\caption{Approximate DP for the random-variance Bayes-optimal threshold.}
\label{alg:adp-random}
\begin{algorithmic}[1]
\State \textbf{Precompute.} Per-stage $\bar X$ grids $\{\bar X_t^{(k_1)}\}$ on $I_{\bar X, t} = [-5 s_{\bar X, t}, +5 s_{\bar X, t}]$ with $s_{\bar X, t} = \sqrt{\tau_0^2 + \sigma_{\max}^2 / t}$; stage-invariant $L$ grid $\{L^{(k_2)}\}$ on $I_L$.
\State \textbf{Initialize.} For each grid point $(k_1, k_2)$, recover $(S_{n-1}, Q_{n-1})$ from $(t = n-1, \bar X_{n-1}^{(k_1)}, L^{(k_2)})$, compute $p(\sigma \mid X_{1:n-1})$ via~\eqref{eq:marginal-loglik}, set $\widehat C_{n-1}(\bar X_{n-1}^{(k_1)}, L^{(k_2)}) \gets \int p(\sigma \mid X_{1:n-1})\, \mu_{n-1}^{(\sigma)}\, d\sigma$
\For{$t = n-2, n-3, \ldots, 1$}
  \For{each grid point $(k_1, k_2)$ on stage $t$}
    \State Recover $(S_t, Q_t)$ from $(t, \bar X_t^{(k_1)}, L^{(k_2)})$
    \State Compute mixture weights $\{\tilde w_k\}_{k=1}^M$ from $p(\sigma \mid X_{1:t})$
    \State Compute component means/variances $\{(\mu_t^{(\sigma_k)}, v_t^{(\sigma_k)})\}_{k=1}^M$ via~\eqref{eq:mu-static}
    \State $S \gets 0$
    \For{$k = 1, \ldots, M$, $j = 1, \ldots, J$}
      \State $x_{kj} \gets \mu_t^{(\sigma_k)} + \sqrt{2 v_t^{(\sigma_k)}}\, z_j$
      \State $S_{t+1} \gets S_t + x_{kj}$;\quad $Q_{t+1} \gets Q_t + x_{kj}^2$
      \State Translate to $(\bar X_{t+1}, L_{t+1})$, bilinearly interpolate $\widehat C_{t+1}^{\mathrm{bilin}}$ on stage-$(t{+}1)$'s $\bar X$ grid, clipping at grid boundary
      \State $S \gets S + \tilde w_k \cdot (w_j/\sqrt\pi) \cdot \max\!\big\{x_{kj}, \widehat C_{t+1}^{\mathrm{bilin}}\big\}$
    \EndFor
    \State $\widehat C_t(\bar X^{(k_1)}, L^{(k_2)}) \gets S$
  \EndFor
\EndFor
\State \Return $\{\widehat C_t\}$
\end{algorithmic}
\end{algorithm}

\paragraph{Convergence check.} Solve at $(K_1, K_2, J) = (256, 256, 64)$ and at reference $(512, 512, 128)$; report (i) max absolute disagreement on $\widehat C$ over the grid, and (ii) action-label disagreement on $10^4$ test sequences per regime. Fix the labeling resolution after the operational disagreement falls below $10^{-3}$.

\subsection{Labeling}
\label{sec:labeling}

For each training distribution we precompute the corresponding ADP threshold table $\{\widehat C_t\}$ once: Algorithm~\ref{alg:adp-static} for $\mathcal{D}_1, \mathcal{D}_2, \mathcal{D}_3$ (1D in $\mu_t$, linear interpolation), Algorithm~\ref{alg:adp-random} for $\mathcal{D}_{\mathrm{disc}}, \mathcal{D}_{\mathrm{logu}}$ (2D in $(\bar X_t, L_t)$, bilinear interpolation). Both tables expose the same query: given the belief state at time $t$, return $\widehat C_t$.

For each sequence $X_{1:n}$ from training distribution $\mathcal{D}$, at each $t \in \{1, \ldots, n-1\}$ we update the belief state and query the table for $\widehat C_t$, then emit
\[
y_t^{\mathrm{cv}} \;:=\; \widehat C_t,
\qquad
y_t^{\mathrm{act}} \;:=\; \mathbf{1}[X_t \ge \widehat C_t].
\]
At $t = n$, $y_n^{\mathrm{act}} = 1$ deterministically and $y_n^{\mathrm{cv}} = 0$.

\section{Portfolio of Baseline Algorithms}
\label{sec:baselines}

We report several decision rules: oracular baselines that assume the test-regime $\sigma$ is known, data-only baselines that estimate $\sigma$ from $X_{1:t}$ (directly comparable to the trained model), and scale-free baselines.

\subsection{Known-$\mu$ and known-$\sigma$ thresholds}
\label{sec:known-mu-sigma}

If both $\mu$ and $\sigma$ were known, the offers $X_1, \ldots, X_n$ would be uninformative about $\mu$ and the optimal continuation value $c_t := \mathbb{E}[\max\{X_{t+1}, c_{t+1}\}]$ would be a deterministic sequence with $c_{n-1} = \mu$. To evaluate the recursion, $\mathbb{E}[\max\{X, a\}] = \mu + \sigma\psi((a - \mu)/\sigma)$ for $X \sim \mathcal{N}(\mu, \sigma^2)$, where $\psi(z) := z\Phi(z) + \varphi(z)$. The standardized continuation value $\eta_t := (c_t - \mu)/\sigma$ depends on neither $\mu$ nor $\sigma$:
\begin{equation}
\label{eq:eta}
\eta_{n-1} = 0, \qquad \eta_t = \psi(\eta_{t+1}), \qquad t = n-2, \ldots, 1.
\end{equation}
The known-$\mu$, known-$\sigma$ optimal rule is $\mathbf{1}[X_t \ge \mu + \sigma\eta_t]$. Baselines below plug different estimators of $\mu$ and $\sigma$ into this rule.

\subsection{Known-$\sigma$ baselines}

Evaluated using the true test-regime $\sigma$.

\paragraph{Plug-in.} $\pi_t^{\mathrm{plug\text{-}in}}(X_{1:t}; \sigma) = \mathbf{1}[X_t \ge \bar X_t + \sigma\eta_t]$. Replaces $\mu$ by the sample mean $\bar X_t$.

\paragraph{Prior-only.} $\pi_t^{\mathrm{prior}}(X_t; \sigma) = \mathbf{1}[X_t \ge \mu_0 + \sigma\eta_t]$. Replaces $\mu$ by the prior mean.

\paragraph{Myopic.} $\pi_t^{\mathrm{myopic}}(X_{1:t}; \sigma) = \mathbf{1}[X_t \ge \mu_t]$, with $\mu_t$ the conjugate posterior mean at known $\sigma$.

\paragraph{Per-regime Bayes-optimal oracle $\pi^\star_{\mathcal{D}_i}$.} The Bayes-optimal policy from Algorithm~\ref{alg:adp-static} at $\sigma = \sigma_i$. The highest expected payoff achievable on $\mathcal{D}_i$ by any online policy.

\subsection{Data-only baselines}

These see only $X_{1:t}$ and estimate $\sigma$ from data --- directly comparable to the trained model.

\paragraph{MAP-$\sigma$ plug-in.} Compute the posterior $p(\sigma \mid X_{1:t})$ via~\eqref{eq:marginal-loglik}, take its mode --- the maximum-a-posteriori estimate
\[
\hat\sigma_t^{\mathrm{MAP}} \;:=\; \arg\max_\sigma p(\sigma \mid X_{1:t}),
\]
and apply
\[
\pi_t^{\mathrm{plug\text{-}in,\,MAP}}(X_{1:t}) \;=\; \mathbf{1}[X_t \ge \bar X_t + \hat\sigma_t^{\mathrm{MAP}}\, \eta_t].
\]

\paragraph{MLE-$\sigma$ plug-in.} The maximum-likelihood estimate of $\sigma$ is the within-sequence sample standard deviation
\[
\hat\sigma_t^{\mathrm{MLE}} \;:=\; \sqrt{\hat\sigma_t^2} \;=\; \sqrt{\frac{1}{t-1}\!\left(Q_t - t\bar X_t^2\right)}.
\]
Apply
\[
\pi_t^{\mathrm{plug\text{-}in,\,MLE}}(X_{1:t}) \;=\; \mathbf{1}[X_t \ge \bar X_t + \hat\sigma_t^{\mathrm{MLE}}\, \eta_t].
\]

\paragraph{Random-variance Bayes-optimal oracle $\pi^\star_{\mathcal{D}}$.} The Bayes-optimal policy from Algorithm~\ref{alg:adp-random} for $\mathcal{D} \in \{\mathcal{D}_{\mathrm{disc}}, \mathcal{D}_{\mathrm{logu}}\}$. Strict upper bound on every online baseline that does not know $\sigma$, under the prior of training distribution $\mathcal{D}$.

\subsection{Scale-free baselines}

\paragraph{Secretary.} Skip the first $r := \lfloor n/e \rfloor$ observations, set $M_r := \max_{s \le r} X_s$, accept the first $X_t$ ($t > r$) with $X_t > M_r$; if none, accept $X_n$.

\paragraph{Offline (hindsight).} $\pi^{\mathrm{offline}}(X_{1:n}) = \arg\max_t X_t$. Not implementable online; $\mathbb{E}[\max_t X_t]$ upper-bounds every online policy. By scale-and-shift, $\max_t X_t = \mu + \sigma Z_n$ where $Z_n$ is the maximum of $n$ iid $\mathcal{N}(0, 1)$ variables, and since $\mathbb{E}[\mu] = \mu_0 = 0$,
\[
\mathbb{E}[\max_t X_t;\, \mathcal{D}_i] \;=\; \sigma_i \cdot \mathbb{E}[Z_n], \qquad \mathbb{E}[Z_n] \;=\; \int_{-\infty}^{\infty} z \cdot n\Phi(z)^{n-1}\varphi(z)\, dz.
\]
The integral has no elementary closed form but is evaluated to machine precision by 1D quadrature.

\section{Model and Training}
\label{sec:training}

\subsection{Architecture}

GPT-2 decoder backbone with $L = 8$ layers, embedding dimension $d_{\mathrm{emb}} = 128$, $M = 4$ heads, causal attention mask, no dropout. Roughly $1.6$M parameters. The standard GPT-2 language-modeling head is removed; the model is not trained to predict $X_{t+1}$ from $X_{1:t}$. Each scalar $X_t$ is mapped to the embedding by a learnable affine layer $\mathbf{e}_t = W_{\mathrm{in}}\, X_t + \mathbf{b}_{\mathrm{in}}$, followed by learned positional embeddings. The hidden state $\mathbf{h}_t$ at each position feeds either a continuation-value head $\widehat C_t = \mathbf{w}_{\mathrm{cv}}^\top \mathbf{h}_t + b_{\mathrm{cv}}$ or an action head $\hat a_t = \sigma(\mathbf{w}_{\mathrm{act}}^\top \mathbf{h}_t + b_{\mathrm{act}})$, trained against the oracle labels (Section~\ref{sec:labeling}). The causal mask enforces that the output at position $t$ depends only on $X_{1:t}$, matching the online structure of the stopping problem. Inputs $X_t$ enter raw; the cv head produces $\widehat C_t$ in raw oracle units. No input-scale normalization is applied.

\subsection{Training objective}
\label{sec:training-objective}

We train ten models: five training distributions $\{\mathcal{D}_1, \mathcal{D}_2, \mathcal{D}_3, \mathcal{D}_{\mathrm{disc}}, \mathcal{D}_{\mathrm{logu}}\}$ $\times$ two supervisions $\{\text{cv}, \text{act}\}$. The model emits one prediction per timestep $t$ via the causal mask. We supervise at every $t \in \{1, \ldots, n-1\}$ (the terminal $t = n$ has no decision, since the agent must accept $X_n$): each timestep contributes its own oracle label and its own gradient signal, giving $n-1 = 255$ training signals per sequence. The per-sequence losses, with sequence-specific $\sigma_i$ used only for cv-loss normalization (the model never sees $\sigma_i$), are
\begin{align*}
\ell_{\mathrm{cv}}(\theta;\, X_{1:n},\, \sigma_i)
&\;=\; \frac{1}{(n-1)\,\sigma_i}\sum_{t=1}^{n-1}\!\Big(\widehat C_t(X_{1:t};\theta) - y_t^{\mathrm{cv}}\Big)^{\!2}, \\
\ell_{\mathrm{act}}(\theta;\, X_{1:n})
&\;=\; -\frac{1}{n-1}\sum_{t=1}^{n-1}\!\Big[\, y_t^{\mathrm{act}}\,\log \hat a_t + (1 - y_t^{\mathrm{act}})\,\log(1 - \hat a_t)\,\Big].
\end{align*}
Each is the within-sequence mean over per-step terms. The mini-batch loss is the further mean over $B = 64$ sequences sampled fresh from the training distribution.

Without normalization, cv squared errors have typical magnitude $\sim \sigma^2$, so a $\sigma = 100$ batch contributes gradients $10^4 \times$ larger than a $\sigma = 1$ batch and a single Adam learning rate cannot accommodate both. We divide by $\sigma_i$ rather than $\sigma_i^2$ to make per-sequence gradient magnitudes regime-invariant: at any step the residual $\widehat C_t - y_t^{\mathrm{cv}}$ has typical scale $\sigma_i$, so $d\ell_{\mathrm{cv}}/d\theta \propto (1/\sigma_i)\cdot\text{residual}\cdot d\widehat C_t/d\theta$ is order unity in $\sigma_i$.\footnote{$1/\sigma_i$ instead of $1/\sigma_i^2$ gives gradient invariance across $\sigma$ regimes; loss values now scale as $\sigma$ and are not directly comparable across regimes, but loss-value comparability across regimes was already not held under $1/\sigma_i^2$ either (only each model's own descent shape is meaningful as a convergence proxy). Logged in the V3 README's spec corrections.} The normalizer is per-sequence (constant within a sequence) and is applied on the loss side only: it gives the model no in-sequence information about $\sigma_i$ --- the model still has to recover $\sigma$ from $X_{1:t}$ on its own. Act supervision is naturally regime-invariant. One consequence: cv val MSE values are in $\sigma^{-1}$ units and not comparable across regimes; only each model's own MSE descent indicates convergence.

Adam, peak learning rate $10^{-4}$, no weight decay, batch size $64$. Static-variance cv models train for $2 \times 10^5$ steps; random-variance cv models for $3 \times 10^5$ (additional steps for in-context regime estimation). Act models train for half their cv counterpart's step count. Cosine schedule with $20\%$ linear warmup; learned absolute position embeddings; checkpoint selection by minimum validation loss.

\section{In-Context Learning: What and How?}
\label{sec:experiments}

All ten models are evaluated on each of the three test regimes $\mathcal{D}_1, \mathcal{D}_2, \mathcal{D}_3$. We split the analysis into two experiments: the six static-variance models (Section~\ref{sec:exp-static}) and the four random-variance models (Section~\ref{sec:exp-random}).

\subsection{Metrics}

For each (model, test regime) pair we report expected payoff $R(\pi; \mathcal{D}_i) := \mathbb{E}[X_{\tau_\pi}]$, normalized payoff $R/R^\star_{\mathcal{D}_i}$ relative to the per-regime oracle (the regime-invariant metric for cross-regime comparison), and per-step action agreement with each baseline. We report agreement as a diagnostic --- two policies that reject most positions can agree on $> 98\%$ of steps while differing on the few accepts that determine the entire payoff, so high agreement does not imply matched payoff --- but $R/R^\star$ is the headline metric.

\subsection{Static-variance experiment}
\label{sec:exp-static}

The six static-variance models test whether single-regime training induces a regime-invariant Bayesian algorithm or merely a fixed-scale policy specific to the training regime. Each model is trained on one of $\mathcal{D}_1, \mathcal{D}_2, \mathcal{D}_3$ and evaluated on all three test regimes, with all baselines reported per (model, test-regime) cell. The off-diagonal $R/R^\star$ values are the key diagnostic.

\subsection{Random-variance experiment}
\label{sec:exp-random}

The four random-variance models test whether training on a non-degenerate $\sigma$-prior induces in-context regime inference. Each model is trained on one of $\mathcal{D}_{\mathrm{disc}}, \mathcal{D}_{\mathrm{logu}}$ and evaluated on all three test regimes, with all baselines reported per cell.

\paragraph{Comparison to data-only baselines.} The MAP-$\sigma$ and MLE-$\sigma$ plug-in baselines see exactly the same input as the model: $X_{1:t}$. The MAP-$\sigma$ baseline uses the same prior over $\sigma$ as the model's training distribution. The trained model can outperform these baselines either by being closer to its training-distribution oracle $\pi^\star_{\mathcal{D}}$, or by exploiting features of $X_{1:t}$ beyond $(\bar X_t, \hat\sigma_t^2)$.

\subsection{Threshold trajectories}
\label{sec:exp-thresh}

For each cv-supervised model and each test regime, plot the per-step threshold (mean $\pm 1$ std over the $N_{\mathrm{test}}$ test sequences) against the per-regime oracle, the random-variance oracle (where applicable), the data-only baselines (MAP-$\sigma$, MLE-$\sigma$), the known-$\sigma$ baselines, the secretary baseline (the running max $M_r$ over the first $r = \lfloor n/e \rfloor$ observations, treated as $-\infty$ for $t \le r$ and held constant for $t > r$), and the offline-hindsight $\mathbb{E}[\max_{1 \le s \le n} X_s; \mathcal{D}_i]$ as a horizontal line per regime (computed via Section~\ref{sec:baselines}). A model that has learned in-context $\sigma$ inference produces threshold trajectories that adapt to the test regime within a few steps; one that has not produces a flat regime-invariant threshold.

\subsection{Attention analysis}
\label{sec:exp-attn}

The diagnostic question is whether the model develops heads that compute $\hat\sigma_t^2$. This requires multiplicative dependence on $X_s$ (squaring), distinct from the uniform-attention running-sum head needed to compute $\bar X_t$. At each position $t$, the causal attention mechanism in layer $\ell$, head $m$ produces a probability distribution $\alpha_t^{(\ell,m)} \in \Delta^{t-1}$ over the past positions $s \le t$. The number of past positions grows with $t$, so direct cross-$t$ averaging conflates structure with horizon length; we instead summarize each $\alpha_t^{(\ell,m)}$ by its normalized Shannon entropy
\[
\widetilde H_t^{(\ell, m)} \;:=\; \frac{H(\alpha_t^{(\ell, m)})}{\log(t+1)} \;\in\; [0, 1],
\]
and plot $\widetilde H_t^{(\ell, m)}$ averaged over the $N_{\mathrm{test}}$ test sequences as a function of $t$. The normalization makes the score horizon-invariant: a uniform head has $\widetilde H_t \approx 1$ at every $t$, a head concentrating on a single past position has $\widetilde H_t \approx 0$. A flat curve at $\widetilde H_t \approx 1$ identifies a candidate running-mean head; structured or low-entropy curves identify candidates for non-trivial computation. Whether such a head computes $\hat\sigma_t^2$ is then tested by probing (Section~\ref{sec:exp-probe}).

\subsection{Probing analysis}
\label{sec:exp-probe}

We freeze a trained model and train auxiliary probe networks to recover target quantities from its hidden states $H^{(\ell)} \in \mathbb{R}^{n \times d_{\mathrm{emb}}}$ at each layer $\ell$.

\paragraph{Probe.} For target $v$ (a scalar or low-dimensional vector that depends on $X_{1:t}$), the probe is
\[
\alpha = \mathrm{softmax}(s), \qquad \hat v = \mathrm{FF}(\alpha^\top W_v H^{(\ell)}),
\]
with learned position-attention weights $s \in \mathbb{R}^n$, projection $W_v \in \mathbb{R}^{n \times d_{\mathrm{emb}}}$, and feed-forward head $\mathrm{FF}$. We train two variants per layer: $\mathrm{FF}$ a linear map (testing whether the target is encoded linearly) and $\mathrm{FF}$ a 2-layer MLP with hidden width $512$ and GeLU activation (testing whether the target is encoded at all). Loss is $\|v - \hat v\|^2$; optimizer is Adam at learning rate $10^{-3}$. The probe's MSE quantifies encoding strength; the position-attention $\alpha$ identifies the timestep whose hidden state carries the target.

\paragraph{Targets.} For each model we probe three tiers of quantities, each computable from $X_{1:t}$ and known oracle parameters:
\begin{itemize}
\item \textbf{Sufficient statistics:} $S_t$, $\bar X_t$, and (for random-variance models) $Q_t$, $\hat\sigma_t^2$. $\hat\sigma_t^2$ requires squaring inputs, so a successful MLP probe alongside a failing linear probe is direct evidence of nonlinear computation.
\item \textbf{Posterior and point estimates of $\sigma$:} $\mu_t$, $\log p(\sigma_i \mid X_{1:t})$ for $i \in \{1, 2, 3\}$, $\hat\sigma_t^{\mathrm{MAP}}$, and $\hat\sigma_t^{\mathrm{MLE}} = \sqrt{\hat\sigma_t^2}$. Probing MAP and MLE jointly discriminates how the model treats $\sigma$: only MAP decodes if the model used the prior; only MLE if it computed the sample SD without a Bayesian update; both, if either is recoverable from a common intermediate.
\item \textbf{Continuation value:} the oracle threshold $\widehat C_t$, plus the data-only baselines' thresholds $\widehat C_t^{\mathrm{plug\text{-}in,\,MAP}} := \bar X_t + \hat\sigma_t^{\mathrm{MAP}}\,\eta_t$ and $\widehat C_t^{\mathrm{plug\text{-}in,\,MLE}} := \bar X_t + \hat\sigma_t^{\mathrm{MLE}}\,\eta_t$, to test whether the internal threshold matches the oracle or a simpler plug-in. Probing the act-supervised models is particularly informative since $\widehat C_t$ never appears in their training signal.
\end{itemize}

\paragraph{Control task.} Probe success on a deterministic function of the input is partly a measure of probe capacity. We  train a control model on a degenerate task where the target threshold is constant in $X_{1:t}$; probes trained on it establish the noise floor that the trained-model probe must beat.

\paragraph{Reporting.} Per (model, target, layer): linear- and MLP-probe MSE on $N_{\mathrm{test}}$ held-out sequences, the position-attention $\alpha$ averaged across sequences, and the control-task MSE.

\subsection{Causal interventions}
\label{sec:exp-intervention}

Probing is correlational: a probe direction does not establish that the model \emph{uses} the encoded quantity to set its threshold. We test causal use via vector arithmetic on the residual stream.

For a target $v$ (e.g., $\hat\sigma_t^2$ or $\hat\sigma_t^{\mathrm{MAP}}$) with linear probe direction $\mathbf{w}$ at layer $\ell$, modify the residual stream during the forward pass at every layer $\ell' \in \{\ell, \ldots, L\}$:
\[
\mathbf{h}_t^{(\ell')} \;\leftarrow\; \mathbf{h}_t^{(\ell')} + \alpha\, \mathbf{w}.
\]
Applying at multiple layers (rather than just $\ell$) is necessary because subsequent layers can otherwise undo the intervention. Choose the scale $\alpha^\star$ so that pushing along $\mathbf{w}$ moves the probe readout from the typical $\hat\sigma_t^2$ at $\sigma = 10$ to its typical value at $\sigma = 100$.

\paragraph{Diagnostic.} For each random-variance model, compute three threshold trajectories averaged over the test set: $\widehat C_t^{\sigma = 10}$ on $\mathcal{D}_2$ sequences (no intervention), $\widehat C_t^{\mathrm{push}}$ on the same $\mathcal{D}_2$ sequences with $\alpha = \alpha^\star$, and $\widehat C_t^{\sigma = 100}$ on $\mathcal{D}_3$ sequences (no intervention). Plot all three, and report the relative shift
\[
\rho_t \;:=\; \frac{\widehat C_t^{\mathrm{push}} - \widehat C_t^{\sigma = 10}}{\widehat C_t^{\sigma = 100} - \widehat C_t^{\sigma = 10}}.
\]
$\rho_t = 0$ means the intervention had no effect; $\rho_t = 1$ means the model's threshold moved entirely from the $\sigma = 10$ regime to the $\sigma = 100$ regime, a signature of a causally-used variance estimate.

\subsection{Logit lens for the continuation value head}
\label{sec:exp-logit-lens}

The probing analysis tells us how target quantities are encoded layer-by-layer; ``the logit lens'' tells us when the model's \emph{output} is finalized. We apply the trained cv head $\mathbf{w}_{\mathrm{cv}}^\top \mathbf{h} + b_{\mathrm{cv}}$ to the residual stream at every intermediate layer, not just the last:
\[
\widehat C_t^{(\ell)} \;:=\; \mathbf{w}_{\mathrm{cv}}^\top \mathbf{h}_t^{(\ell)} + b_{\mathrm{cv}}, \qquad \ell = 0, 1, \ldots, L.
\]
The trajectory $\ell \mapsto \widehat C_t^{(\ell)}$ shows how the predicted threshold evolves through depth. For act-supervised models we apply the act head analogously to recover $\hat a_t^{(\ell)}$.

\paragraph{Diagnostic.} Per (model, test regime), plot the per-step trajectory of $\widehat C_t^{(\ell)}$ against $\ell$, averaged over the test set. The earliest layer at which $\widehat C_t^{(\ell)}$ matches the final-layer output to within a fixed tolerance (e.g., $1\%$) localizes the depth at which the model commits to its threshold. Combined with the layerwise probe-MSE curves of Section~\ref{sec:exp-probe}, this orders the depths at which sufficient statistics, posterior estimates of $\sigma$, and the threshold itself become decodable.

\subsection{Activation patching}
\label{sec:exp-patching}

The residual-stream interventions of Section~\ref{sec:exp-intervention} push along a probe direction; activation patching localizes the same causal claim to specific components. Run the model on a clean sequence drawn from $\mathcal{D}_2$ (true noise level $\sigma = 10$) and a corrupt sequence drawn from $\mathcal{D}_3$ (true $\sigma = 100$). For each candidate patch site --- a (layer $\ell$, position $t'$, component $c$) triple where $c \in \{\text{attention}, \text{MLP}, \text{head}\;m\}$ --- replace the clean run's activation at that site with the corresponding corrupt activation, then complete the forward pass on the clean sequence and record the resulting $\widehat C_t^{\mathrm{patch}}$.

We report the relative shift
\[
\rho_t^{\mathrm{patch}}(\ell, t', c) \;:=\; \frac{\widehat C_t^{\mathrm{patch}} - \widehat C_t^{\mathrm{clean}}}{\widehat C_t^{\mathrm{corrupt}} - \widehat C_t^{\mathrm{clean}}}
\]
averaged over the test set, on the same $[0, 1]$ scale as Section~\ref{sec:exp-intervention}. Sites with $\rho_t^{\mathrm{patch}} \approx 1$ are individually sufficient to transport the threshold from the $\sigma = 10$ regime to $\sigma = 100$. Heatmaps over $(\ell, t')$ for each component class identify the heads or MLPs that carry the noise-level signal --- a sharper localization than the residual-stream pushing, which acts on all components at once.

\subsection{Head ablation at inference time}
\label{sec:exp-ablation}

Probing tells us what is encoded; ablation tells us what is used. Take a trained random-variance model and, at inference time on a held-out test set, zero or mean-replace the output of one attention head $(\ell, m)$ at all positions, leaving every other head untouched. Measure the resulting $R/R^\star$ on each test regime, and report the per-head, per-regime performance drop $\Delta_{\ell,m}^{(i)} := R/R^\star\big|_{\text{full}} - R/R^\star\big|_{\text{ablate}(\ell, m)}$ on $\mathcal{D}_i$.

The diagnostic is not just \emph{which heads are necessary} but \emph{for which regime}. A head whose ablation drops performance uniformly across $\mathcal{D}_1, \mathcal{D}_2, \mathcal{D}_3$ is doing regime-invariant work (e.g., computing $\bar X_t$). A head whose ablation drops performance only on $\mathcal{D}_3$, or whose drop scales with $\sigma$, is specifically a noise-level-handling head --- evidence that the model has learned regime-specialized circuitry. Because this requires no retraining, it can be run cheaply across all $L \times M = 32$ heads.

\appendix

\section{Derivation of the marginal log-likelihood}
\label{app:marginal-derivation}

We derive equation~\eqref{eq:marginal-loglik} by simplifying the determinant and inverse of the covariance $\Sigma_t^{(\sigma)} = \sigma^2 I_t + \tau_0^2 \mathbf{1}\mathbf{1}^\top$. Since $\Sigma_t^{(\sigma)}$ is a rank-one update of a scaled identity, both the matrix determinant lemma and the Sherman--Morrison formula apply: for invertible $A$ and column vectors $u, v$,
\[
\det(A + uv^\top) \;=\; \det(A) \cdot (1 + v^\top A^{-1} u),
\qquad
(A + uv^\top)^{-1} \;=\; A^{-1} - \frac{A^{-1} u v^\top A^{-1}}{1 + v^\top A^{-1} u}.
\]
In our case $A = \sigma^2 I_t$, $u = v = \tau_0 \mathbf{1}$, so $u v^\top = \tau_0^2 \mathbf{1}\mathbf{1}^\top$ and $\Sigma_t^{(\sigma)} = A + uv^\top$. The starting point is the log density of $X_{1:t} \mid \sigma \sim \mathcal{N}(\mathbf{0}, \Sigma_t^{(\sigma)})$,
\[
\log p(X_{1:t} \mid \sigma) \;=\; -\frac{t}{2}\log(2\pi) - \frac{1}{2}\log\det\Sigma_t^{(\sigma)} - \frac{1}{2} X_{1:t}^\top\!\big(\Sigma_t^{(\sigma)}\big)^{\!-1} X_{1:t}.
\]

For the determinant: $A^{-1} = \sigma^{-2} I_t$ gives $v^\top A^{-1} u = \tau_0^2 \mathbf{1}^\top \mathbf{1}/\sigma^2 = t\tau_0^2/\sigma^2$, and $\det A = \sigma^{2t}$, so by the matrix determinant lemma
\[
\log\det\Sigma_t^{(\sigma)} \;=\; t\log\sigma^2 + \log\!\left(1 + \frac{t\tau_0^2}{\sigma^2}\right).
\]
For the inverse: $A^{-1} u = \sigma^{-2}\tau_0\mathbf{1}$ and $A^{-1} u v^\top A^{-1} = \tau_0^2 \mathbf{1}\mathbf{1}^\top/\sigma^4$, so by Sherman--Morrison
\[
\big(\Sigma_t^{(\sigma)}\big)^{\!-1} \;=\; \frac{1}{\sigma^2}I_t - \frac{\tau_0^2/\sigma^4}{1 + t\tau_0^2/\sigma^2}\mathbf{1}\mathbf{1}^\top.
\]
The quadratic form follows from $X_{1:t}^\top I_t X_{1:t} = Q_t$ and $X_{1:t}^\top \mathbf{1}\mathbf{1}^\top X_{1:t} = (\mathbf{1}^\top X_{1:t})^2 = S_t^2$:
\[
X_{1:t}^\top\!\big(\Sigma_t^{(\sigma)}\big)^{\!-1}\! X_{1:t} \;=\; \frac{Q_t}{\sigma^2} - \frac{\tau_0^2 S_t^2/\sigma^4}{1 + t\tau_0^2/\sigma^2}.
\]
Assembling the three terms yields
\[
\log p(X_{1:t} \mid \sigma) \;=\; -\frac{t}{2}\log(2\pi) - \frac{t}{2}\log\sigma^2 - \frac{1}{2}\log\!\left(1 + \frac{t\tau_0^2}{\sigma^2}\right) - \frac{1}{2}\!\left[\frac{Q_t}{\sigma^2} - \frac{\tau_0^2 S_t^2/\sigma^4}{1 + t\tau_0^2/\sigma^2}\right],
\]
where the first term is absorbed into $\mathrm{const}(t)$ in equation~\eqref{eq:marginal-loglik}. The dependence on $X_{1:t}$ enters only through $(S_t, Q_t)$.

\bibliographystyle{plainnat}
\bibliography{refs}

\end{document}